\cotefile{156-5}
\scanpage{34}\mspage{33}
\diagram{\begin{tikzcd}[ampersand replacement=\&,column sep=small]
I' \arrow[d,"\varphi"',"\text{surj.}"] \& \hookleftarrow \& J'=\varphi^{-1}(J)\arrow[d] \& \& I'^{*}\arrow[d,"{\varphi^{*}=\varphi\mid I'^{*}}"]\\
I \& \hookleftarrow \& J\ \text{partie f.} \& \text{d'où} \& I^{*} \& \begin{gathered}\text{\deleted{\unsure{croiss.|surj.}}}\\ \text{surjectif}\end{gathered}
\end{tikzcd}}

et on a
\diagram{\begin{tikzcd}[ampersand replacement=\&,column sep=large]
U\arrow[r,"p'","\text{surj}"']\arrow[rr,bend right=23,"p" description,"\text{surj}"'] \& I'^{*}\arrow[r,"\varphi^{*}","\text{surj}"'] \& I^{*}
\end{tikzcd}}
% slanted
\margin{Attention, $\varphi^{*}$ bijectif n'implique pas néc. $\varphi^{*}$ iso !!! Mais $=$ on suppose $\varphi^{*}$ iso, ça n'implique pas a), cf plus bas.}
les \deleted{$A$} strates ouvertes de $F'$ qui sont \deleted{dans} \inserted{$\subset$} $U$ (i.e. pour \deleted{dans} $\in G'$) sont les fibres de $U\xrightarrow{p'}\unsure{I^{*\alpha}|I^{**}|I'^{*}}$, les strates ouvertes de $F$ qui sont $\subset U$ (i.e. pour $\in G$) sont les fibres de $U\xrightarrow{p}I^{*}$.

La condition b) dit que $\varphi^{*}$ est bijective,
La condition c) dit que les \deleted{\unsure{s}}fibres de $p$ (ce sont les $A^\circ$ pour $A\in F\setminus G$) sont des fibres pour $p'$. C'est \deleted{éq} visiblement équivalent.

\deleted{D'autre part,}

D'autre part, que dit d) ? Notons \inserted{(on l'a \unsure{pr. entre|prouvé} pour $F'\ll F$)\deleted{$A\in F$})} que pour $A\in F$\deleted{\illegible{1}}, on a $A\in F'$ sss $A^{\circ(F)}$ est une strate ouverte de $F'$ (cf lemme déjà prouvé, \inserted{qu'on donnera} plus bas). Donc d) dit que pour $\forall\,A_i\in F\setminus G=I^{*}$, \deleted{\unsure{j}} la \inserted{$F$-}strate $p^{-1}(\{\deleted{\unsure{A}}i\})$\inserted{$\subset U$} est une \inserted{$F'$-}strate \deleted{de $p$} de $U$. Donc cela \unsure{équiv. à visiblement à|équiv. visiblement à} $p$ bijectif.

Donc on a
\[b\Longleftrightarrow c\Longleftrightarrow d\]

Montrons que c'est équiv. à a. Soit

\deleted{$F''$ comme ds a), d'où $I''$, et sur} \deleted{\unsure{$S=|F$}}

\deleted{$S=|F|=|F'|=|F''|$, on a}

\diagram{\begin{tikzcd}[ampersand replacement=\&,column sep=large,execute at end picture={\draw (\tikzcdmatrixname-1-1.north west) -- (\tikzcdmatrixname-2-1.south west); \draw (\tikzcdmatrixname-1-2.north west) -- (\tikzcdmatrixname-2-2.south east); \draw (\tikzcdmatrixname-1-3.north east) -- (\tikzcdmatrixname-2-3.south east);}]
\deleted{S}\arrow[r]\arrow[d]\arrow[rr,bend left=18,"\deleted{\varphi}"] \& \deleted{I''}\arrow[r,"\deleted{\varphi''}"]\arrow[d] \& \deleted{I}\arrow[d]\\
\deleted{|G'|=|G|=T}\arrow[r] \& \deleted{\unsure{J''|J'}}\arrow[r] \& \deleted{J}
\end{tikzcd}}

\scanpage{35}\mspage{34}
Interprétons la donnée d'une $F'\preccurlyeq F$ induisant $G'\preccurlyeq G$, en termes des données initiales, exprimées par un diagramme d'applications (où $S=|F|$, $T=|G|$, $I$ et $J$ les ens \inserted{ord} d'indices pour $F,G$, $J'$ celui pour $G'$)

\diagram{\begin{tikzcd}[ampersand replacement=\&,column sep=huge,row sep=large]
\& \deleted{\varphi}f^{-1}(J)=T \arrow[r,hook,no head]\arrow[d,"{\deleted{\unsure{\psi|\varphi}}\unsure{g_J^{*}|g_J}}"]\arrow[dd,bend right=30,"{\deleted{\varphi}\unsure{f_J^{*}|f_J}=\deleted{\varphi}f\mid T}"'] \& S\arrow[dd,"{ \deleted{\varphi}f\ \text{surj.}}"]\\
(*) \& J'\arrow[d,"{\deleted{\chi}h_J}"] \& {}\\
\& J\arrow[r,hook,no head,"i","\substack{\text{partie}\\\text{fermée de}\\I}"'] \& I
\end{tikzcd}}
\margin{\uline{NB} On part de $f\colon S\mathbin{-}I$ et \deleted{factorise} \inserted{$J\subset I$}, d'où $T=f^{-1}(J)$, $f_J\colon T\longrightarrow J$, et on factorise $f_J$ en $T\mathbin{-}J'\mathbin{-}J$}

La donnée de $G'$ équivaut à celle d'une factorisation de $\varphi'=\varphi\mid T$ en $T\xrightarrow{\deleted{\unsure{\psi'}}\unsure{g_J|g_J^{*}}}J'\xrightarrow{\deleted{\unsure{\chi'}}\unsure{h_J^{*}|h_J}}J$, où $\chi'$ est une \deleted{\unsure{un}} appl. croiss. surjective d'ens ordonnés, et $\psi'$ une application ensembliste surjection. La donnée d'une $F'$ équivaut à une factorisation analogue pour $\varphi$ en
\diagram{\begin{tikzcd}[ampersand replacement=\&,column sep=large]
S\arrow[r,"{\deleted{\unsure{\psi}}g}"] \& I'\arrow[r,"{\deleted{\chi}h}"] \& I \& \substack{g\ \text{surj}\\ h\ \text{croiss. surj.}}
\end{tikzcd}}

Dire que ça induit \deleted{\unsure{un diag}} la \deleted{\unsure{situation}} \inserted{sous-d. donnée $G'$ de $G$} donnée signifie qu'\deleted{\unsure{elle}} \deleted{\unsure{a des d'applications}} \inserted{existe $i'\colon J'\hookrightarrow I'$} tel que les deux carrés commutent
\diagram{\begin{tikzcd}[ampersand replacement=\&,column sep=huge]
T\arrow[r,hook,no head]\arrow[d,"g_J"]\arrow[dd,bend right=30,"f_J"'] \& S\arrow[d,"g"]\\
J'\arrow[r,hook,"i'"]\arrow[d,"h_J"] \& I'\arrow[d,"h"] \& \text{avec }\deleted{\unsure{\xi}}J'=\deleted{\unsure{\chi|k}}h^{-1}(J)\\
J\arrow[r,hook,"i"] \& I
\end{tikzcd}}

\scanpage{36}\mspage{35}
\deleted{Si on \deleted{\unsure{en}} a deux \unsure{comme ça|ensembles sc.} $F'$ et $F''$, correspondant à des ens d'indices $I',I''$, dire que $F''\preccurlyeq F'$ signifie que}

\deleted{Donc} Mais ds une telle situation,

Considérons l'ens ordonné suivant, défini en termes de la situation $(\deleted{\unsure{N}}*)$ ``avant $I'$''
\[
 K=J'\mathbin{\amalg} I\setminus J\qquad\text{en \deleted{tant} qu'ens.}
\]
la relation d'ordre sur $K$ induit celles données sur $J',I\setminus J$, et si $\alpha\in J'$, $\beta\in I-J$, on pose
\[
 \alpha\leq\beta\quad\text{sss}\quad h_J(\alpha)\leq\beta\quad\text{dans }I
\]
\deleted{\unsure{pour} \unsure{$\beta$}.}
donc $(\alpha,\beta)\in K\times K$, on a $\alpha\leq\beta$ sss
% slanted
\margin{expliciter ds une proposition à part}
\[
\left|\begin{aligned}
 &\text{ou bien}\quad \deleted{\unsure{\alpha,\beta|\alpha}}\;(\alpha,\beta)\in\deleted{\unsure{I}}J' \ \text{et}\ \alpha\leq\beta\ \text{ds }J',\\
 &\text{ou bien}\quad (\alpha,\beta)\in I\setminus\deleted{\unsure{\alpha}}J\ \text{et}\ \alpha\leq\beta\ \text{dans }I\setminus J\\
 &\text{ou bien}\quad \alpha\in\deleted{\unsure{I}}J',\ \beta\in I\setminus J\quad\text{et}\quad \deleted{\alpha\leq}h_J(\alpha)\leq\beta
\end{aligned}\right.
\]

Ceci posé, on a, en termes de $J'\xrightarrow{h_J}J\overset{i}{\hookrightarrow}I$ une factorisation canonique du composé
\diagram{\begin{tikzcd}[ampersand replacement=\&,column sep=large]
J'\arrow[r,hook,"{i}"]\arrow[d,"h_J"'] \& \overset{\deleted{\unsure{I''|\text{en}}}}{K}\arrow[d,"h_K"]\\
J\arrow[r,hook,no head,"i"] \& I
\end{tikzcd}}
\margin{$h_K$ est $h_J$ sur $J'$ et l'\unsure{idel|ident.} can. sur $I\setminus J$}

On a fait ce qu'il fallait pour que sur $J'$, $K$ induise la relation de $J'$, que $h_K$ induise un iso d'ens ordonnés $K\setminus J'\xrightarrow{\sim}I\setminus J$, \deleted{et que de plus,} \inserted{cela implique que le} \deleted{\unsure{pour tout le} \unsure{G qui soit}} \inserted{carré est cartésien} ($J'=h_K^{-1}(J)$ avec \unsure{rel.|rels} induite) D'autre part, on voit

\scanpage{37}\mspage{36}
aussitôt que $K$ est objet final, dans la catégorie de tous les carrés commutatifs
\diagram{\begin{tikzcd}[ampersand replacement=\&,column sep=large]
\unsure{{+}J'|J'}\arrow[r,"i'"]\arrow[d,"h_J"] \& I'\arrow[d,"h"]\\
J\arrow[r] \& I
\end{tikzcd}}
\deleted{qui en} \inserted{d'ens ordonnés (ou préord)} qui satisfont la condition précédente. Ça devrait impliquer $(b,c,d)\Longleftrightarrow a)$.

\uline{18 Juin}\quad Dans toutes ces questions, le point de vue des ultr. topologies \inserted{(i.e. top. \uuline{pré}ordonnées)} est peut-être le plus pratique, sans prendre la peine tjrs d'expliciter laborieusement en termes d'ensembles d'indices ordonnés. Une ``figure'' \inserted{de $L$} s'explicite comme un couple $(\deleted{A}S,T)$ d'une partie $S$ de $L$, et d'une ultr.topologie $T$ sur $S$ (topologie où \uline{toute} réunion de fermés est fermée, i.e \uline{toute} intersection d'ouverts est ouverte, i.e. l'intersection des voisinages \inserted{[\unsure{ouverts, si on veut|ouverts si on veut}]} de tout point est un voisinage, ou encore que tt pt a un plus petit voisinage (néc. ouvert). Si $(S,T),(S',T')$ sont comme ça, dire qu'ils sont compatibles signifie que $T,T'$ induisent la $=$ topologie sur $S\cap S'$, \inserted{qui est fermé dans $S'$ (pour $T$) et dans $S'$ (pour $T'$)} et alors y a une unique

\scanpage{38}\mspage{37}
topologie sur $S\cup S'\simeq S\amalg_{S\cap S'}S'$ qui \deleted{indui} induise sur $S,S'$ les topologies données, et pour laquelle $S$ et $S'$ soient fermés -- c'est la topologie finale pour
\diagram{\begin{tikzcd}[ampersand replacement=\&,row sep=0pt,column sep=small]
S.\arrow[r] \& {}\\[-.7ex]
\& \phantom{\longrightarrow}S\deleted{\unsure{\cup}}\amalg S',\\[-.7ex]
S'\arrow[r,no head,end anchor=center] \& \unsure{\cdot|{>}}
\end{tikzcd}}
pour $T$ sur $S$, $T'$ sur $S'$.

[ \uline{NB} Si on a des espaces \unsure{ultr.|ultra}topologiques $(S_i,T_i)_{i\inJ}$ et des applications $S_i\overset{\unsure{f_i|f^{\cdot}|f'}}{\mathbin{-}}\deleted{S}X$, la topologie finale sur $S$ pour $(f_i)$, admettant comme fermés les $A\in X$ tels que $f_i^{-1}(A)$ soit fermé pour $\forall i\in I$, est \inserted{une} ultr-topologie.

\deleted{La relation de préordre associée est la suivante :}

Dans le cas d'espèce, $S\cup S'$ s'identifie, à la somme amalgamée de $S,S'$ suivant le sous-espace fermé commun -- opération topologique bien familière ! ]

Si on a $S\inserted{'}\deleted{\unsure{\supset}}\subset S\deleted{'}$, $(S,T)$ et $(S',T')$ sont compatibles i.e. $(S',T')$ est une sous-figure de $(S,\unsure{\Pi|T})$, sss $S'$ est fermé dans $S$ pour $T$, et $T'$ est la top. induite par $T$. Les sous figures de $(S,T)$ s'identifient donc aux parties fermées de $S$, avec la top. induite.

D'autre part, la relation
\[
 (S',T')\ll(S,T)
\]
signifie que $S'\subset S$, et que $S'\mathbin{\unsure{-|\to}}S$ est continue ($\Longleftrightarrow$ croissante pour les préordres associés \unsure{à|:} $T',T$) La relation

\scanpage{39}\mspage{38}
\[
  (S',T')\preccurlyeq(S,T)
\]
signifie que $S'=S$, et $T'$ plus fine que $T$ i.e. l'application identique $S\xrightarrow{\mathrm{id}}S$ est continue.
% margin note written sideways
\margin{Description de $\bigwedge_i(S_i,T_i)$, c'est $S=\bigcap S_i$ muni de la \inserted{\uuline{ultra}}\uline{topologie} \deleted{ultra} \uline{initiale} sur $S$ pour \inserted{la famille des} $S\hookrightarrow S_i$ (muni de $T_i$ : \deleted{cette ultra}top. pour une famille quelc. d'applications \deleted{\unsure{$f$}} $f_i\colon S\mathbin{-}S_i$}

Revenons à la situation (on reprend les notations $F\ldots$)
\diagram{\begin{tikzcd}[ampersand replacement=\&,column sep=small,row sep=tiny]
G \& {\subset} \& F\\
\rotatebox{90}{$\preccurlyeq$}\\
\deleted{\unsure{F}}G'
\end{tikzcd}}
donc, posant $T=|G|$, $S=|F|$
% margin note written sideways
\margin{de $S$ dans des $S_i$ munis de $T_i$ \unsure{existe} toujours, en \uline{ultr.} \uline{topologie} \uline{engendrée} par les \inserted{fermés} $f_i^{-1}(F_i)$ ($F_i$ fermé ds $T_i$) : on prend les \deleted{réunions} \inserted{intersections} quelc. de réunions quelc. d'ensembles de la forme $f_i^{-1}(F_i)\ldots$}
\diagram{\begin{tikzcd}[ampersand replacement=\&,column sep=small,row sep=tiny]
\deleted{\unsure{\text{on a}}} \& T\arrow[r,hook,no head] \& S\\
\& \tau_0^{+} \& \tau\\
\& \rotatebox{90}{$\leq$}\\
\& \tau'_0
\end{tikzcd}}
$T$ fermé de $(S,\tau)$, $\tau_0$ top. induite par $\tau$, $\tau'_0$ top. plus fine. On considère les \deleted{topologies} $F'\preccurlyeq F$ qui induisent $G'$, i.e. les topologies $\tau'$ sur $S$ telles que $\tau'\leq\tau$ (plus fines que $\tau$ -- il y a \uline{plus} de fermés) et qui \deleted{\unsure{in}}induisent $\tau'_0$.
% margin note written sideways
\margin{correspond à la notion de relation de préordre \deleted{\unsure{engendrée}} \inserted{induite} \deleted{\unsure{de des}} \inserted{images inverses} \unsure{(}relations de préordres) : \deleted{donc ici} $x\leq y\Longleftrightarrow f_i(x)\leq f_i(y)$ \ $\forall i\in I$}
Donc on veut que
\[
\begin{aligned}
 1^\circ\quad &(S,\tau')\longrightarrow(S,\tau)\quad\text{continue},\\
 \text{et}\qquad 2^\circ)\quad &(\deleted{\unsure{S}}T,\tau'_0)\longrightarrow(S,\tau')\quad\text{continue}
\end{aligned}
\]
situation mixte à cheval sur celle des top. initiales et finales \ldots
\diagram{\begin{tikzcd}[ampersand replacement=\&,column sep=large,row sep=tiny]
T\arrow[r,"{i}"] \& S\arrow[r,"\mathrm{id}"] \& S\\
\tau'_0 \& \tau' ? \& \tau
\end{tikzcd}}
\margin{on sait que le composé est continu.}

Donc on veut (i) que les $A\in\unsure{P|\mathcal P|\mathscr P}(S)$ fermés pour $\tau$, \unsure{i.e.} $A\in P_\tau(S)$, soient fermés pour $\tau'$, \deleted{i.e.} et (ii) que pour \uline{$\forall$} $A'$ fermé pour $\tau'$, $i^{-1}(A')=A'\cap T$ soit fermé pour $\tau'_0$. \inserted{Si on veut \unsure{$=$'} $\tau'_0=\tau'\mid T$,} $T$ étant fermé \deleted{\unsure{pou}} pour $\tau$ donc pour $\tau'$, les parties \deleted{\unsure{de}} \inserted{fermées} \unsure{de} $T$ \inserted{pour $\tau'_0$} sont fermées

\scanpage{40}\mspage{39}
pour $ \tau'$. Donc cela impose, comme parties fermées, pour $\tau'$, déjà au moins les
\[
\left\{\begin{aligned}
 A\subset S&\quad\text{fermés pour }\tau\\
 B\subset T&\quad\text{fermés pour }\tau'_0,
\end{aligned}\right.
\]
donc les réunions $A\cup B$. Je dis que ceux-ci sont \inserted{déjà} les fermés d'une ultr.topologie (qui visiblement raffine $\tau$, et qui induit la topologie $\tau'_0$, puisque
\[
 (A\cup B)\cap T=(A\cap T)\cup B
\]
et $A\cap T$ est fermé pour $\tau'_0$, donc on trouve exactement les fermés de $(T,\tau'_0)$).
% slanted
\margin{façon \uline{canonique} d'écrire un tel fermé $X$ de $S$ en termes de $\deleted{\unsure{A}}B=X\cap T$ et $U=X\cap(S\setminus T)$, $U$ un \uline{fermé de $S\setminus T$} (pour $\tau\mid S\setminus T$) condition $\boxed{B\supset\overline U\inserted{\text{adh. pour }\tau\text{, bien sûr}}\cap T}$ et on a $X=B\cup\overline U$. Cette fois correspondance biunivoque entre couples $(B,U)$, satisfaisant les conditions dites, et les $X$. (Dans la représentation précédente, on prend $A=\overline U$, donc on suppose $1^\circ)\ B\supset A\cap T\quad 2^\circ)\ A\cap(S\setminus T)$ dense \unsure{en|dans} $A$.}
\unsure{ce sera|c'est} donc la \deleted{plus} moins fine des topologies cherchées.

\deleted{\unsure{Déjà}} Stabilité par sommes quelconques : triviale. Stabilité par intersections quelconques
\[
 \bigcap_{i\in I}(A_i\cup B_i)=
 \bigcup_{J\in I}
 \biggl(\underset{\underset{A}{\rule{.35pt}{.3em}\rule{3em}{.35pt}\rule{.35pt}{.3em}}}{\bigcap_{i\in J}A_i\biggr)}
 \cap\biggl(\underset{\underset{B}{\rule{.35pt}{.3em}\rule{2.8em}{.35pt}\rule{.35pt}{.3em}}}{\bigcap_{\unsure{j|J}\in I\setminus J} B_J}\biggr)
\]
c'est une réunion de fermés de $(S,\tau)$, donc un fermé de $(S,\tau)$, OK.

Quelle est la topologie induite \inserted{par $\tau'$} sur $S\setminus T$ ? $(A\cup B)\cap(S\setminus T)=A\cap(S\setminus T)$, on trouve la top. induite par $\tau$.

\deleted{Inversement, si $\tau'\leq\tau$, $\tau'\mid T=\tau'_0$, et $\tau'\mid S\setminus T=\tau\mid S\setminus\deleted{\unsure{J}}T$, je dis que $\tau'$ est la top. précédente. \unsure{Il reste à|Je voudrais} voir que ses fermés sont de la forme $A\cup B$. Mais soit un tel fermé $C$, alors \deleted{\unsure{$C=C$}}}

\scanpage{41}\mspage{40}
\deleted{\inserted{car $C$ est $ \tau'$-fermé}}
\[
\begin{aligned}
 &\deleted{\unsure{C'|C}=}\ \underset{\substack{\text{\deleted{fermé de}}\\\deleted{(T,\tau'_0)}\\\text{\deleted{soit }}\deleted{B}}}{\underset{\rule{.35pt}{.35em}\rule{4.4em}{.35pt}\rule{.35pt}{.35em}}{\deleted{(C\cap T)}}}
 \ \deleted{\cup}\ 
 \underset{\substack{\text{\deleted{fermé de} \deleted{\inserted{partie de $S$ dont la trace sur}}}\\
 \text{\deleted{\inserted{$T$ est $\unsure{\tau'_0|\tau_0}$-fermée, et}}}\\
 \text{\deleted{$(S,\tau)$, soit} \deleted{\inserted{la trace sur $S\setminus T$}}}\\
 \text{\deleted{\inserted{est égale à $S\setminus T$}}}\\\deleted{A}}}{\underset{\rule{.35pt}{.35em}\rule{7.8em}{.35pt}\rule{.35pt}{.35em}}{\deleted{\overline{(C\cap(S\setminus T))}^{(\tau')}}}}
\end{aligned}
\]
\deleted{OK. Je voudrais résumer ce qui a été vu dans un lemme topologique plus général :}

\uline{Lemme}\quad \unsure{Soient|Soit} $(\deleted{\unsure{T}}S,\tau)$ espace topologique, $T\deleted{\unsure{_0|_{\varnothing}}}\subset S$ une partie fermée, et $\tau'_0$ une topologie \deleted{\unsure{p}} sur $T$ plus fine que la top. induite (i.e.
\[
 (T,\tau'_0)\mathbin{\longrightarrow}(S,\tau)\qquad\text{continue).}
\]

\deleted{Soit $\tau'$ une topologie sur}

\deleted{Il y a alors \unsure{d'}} Soit \uline{$\mathscr{T}$} l'ens des topologies sur $S$ qui sont

$1^\circ)$ plus fines que $\tau$, et

$2^\circ)$ induisent $\tau'_0$ sur $T$.

\textcircled{a} Il y a alors dans \uline{$\mathscr T$} une topologie la moins fine que toutes, soit $\tau'_{a}$.

\textcircled{b} Pour qu'une partie $X$ de $S$ soit fermée pour $\tau$, il f et s. que (posant $U=S\setminus T$)
\[
\begin{aligned}
1^\circ)\quad &\deleted{\unsure{A}}X\cap \deleted{B}U\deleted{\longrightarrow}\overset{\inserted{\mathrm{dfn}}}{=}X_{u}\quad\text{soit fermé dans }U\quad(\text{pour }\tau\mid U)\\
2^\circ)\quad &X\deleted{_a}\cap\deleted{B}T\overset{\inserted{\mathrm{dfn}}}{=}X_{\deleted{B}T}\ \text{soit fermé dans }T\quad(\text{pour }\tau'_0)\\
3^\circ)\quad &\deleted{\overline{X\cap U}}\ \overline{X_u}^{(\tau)}\deleted{\unsure{\subset|\tau}}{}\cap T\subset X_{\deleted{B}T}.
\end{aligned}
\]
\deleted{\textcircled{c} Les fermés de $\tau'$ sont \unsure{écrits} de façon canonique sous la forme \unsure{$+$|$t$}}

\scanpage{42}\mspage{41}
c) L'ens des fermés pour $\tau'$ est en corr. biuniv. avec l'ens des couples
\[
\begin{aligned}
(B,A_{u}) \qquad & B\subsetT\ \text{fermé pour }\tau'_0\\
& A_u\ \text{fermé de }U=S\setminus T\ \text{pour }\tau\mid(S\setminus T)
\end{aligned}
\]
\[
\text{tels que}\qquad \overline{A_u}\cap T\subset B \qquad
\text{(adh. prise pour }\tau\text{)}\ \deleted{\tau)}.
\]

d) Si $\tau'_0$ et $\tau$ ultra topologiques, $\tau'$ aussi. La relation de préordre associée \deleted{se définit} \inserted{se décrit} ainsi. Pour $x,y\in S$, on a $x\underset{\tau'}{\leq}y$ sss
\[
\left\{\begin{aligned}
&\text{ou bien}\quad x,y\in T,\qquad x\underset{\tau'}{\leq}y\\
&\left.\begin{aligned}
&\text{ou bien}\quad x,y\in U=S\setminus T,\qquad x\underset{\tau}{\leq}y\\
&\text{ou bien}\quad x\in T,\ y\in U,\qquad x\underset{\tau}{\leq}y
\end{aligned}\right]
\end{aligned}\right.
\]
\margin{en d'autres termes, si $y\in U$, on a $x\underset{\tau'}{\leq}y$ sss \unsure{a} $x\underset{\tau}{\leq}y$.}

e) \inserted{Sous les conditions de d)} Considérons les ens ordonnés $I,I'$ associés à $\tau,\tau'$, et la partie \inserted{fermée} $J\subset I$ correspondant à la partie fermée $T$ de $(S,\tau)$, enfin l'application $\varphi:I'\mathbin{\text{---}}I$ canonique.
\diagram{\begin{tikzcd}[ampersand replacement=\&,column sep=small,row sep=large]
\varphi^{-1}(J)={} \& J' \arrow[d] \& \subset \& I' \arrow[d,"\varphi"] \& \overline{\tau'} \arrow[d] \\
{} \& J \& \subset \& I \& \overline{\tau}
\end{tikzcd}}

\deleted{Mais}, \deleted{\unsure{la|le}} Soit $\overline{\tau'}$ \deleted{est} la top de $\tau'$ déduite de $\tau'$ par passage au quotient (ou encore, associée à la relation \uline{d'ordre} sur $I'$ associé au préordre), $\widetilde{\tau}$ la topologie sur $I'$ image inverse de $\overline{\tau}$. Alors $\overline{\tau'}$ est la moins fine des top. sur $I'$ qui \inserted{sont plus fines que $\widetilde{\tau}$ et} induisent sur $J'$ la $=$ \deleted{top} topologie $\overline{\tau}\mid J'$. Elle est

\scanpage{43}\mspage{42}
\uline{caractérisée} par les conditions suivantes :
% slanted
\margin{C'est faux, ça implique (i) (ii) (iii) (iv), mais ces conditions ne suffisent nullement, voir plus bas ! v.~48}
\margin{i.e $\overline{\tau'}\in\underline{\underline{\overline{\mathscr{T}}}}$ $\left\{\vphantom{\rule{0pt}{4ex}}\right.$}

(i) $\overline{\tau'}$ est plus fine que $\widetilde{\tau}$, et induit sur $J'$ $\widetilde{\tau}\mid J'$.

(ii) $\overline{\tau'}\mid\overset{\deleted{\unsure{\sim|\mathord{-}}}}{U'}=I'\setminus J'$ coïncide avec $\widetilde{\tau}\mid\overset{\deleted{\sim}}{U'}$

(iii) $\varphi\mid U'$ induit \deleted{une bijection} \inserted{\unsure{iso. préord.}} $U'\longrightarrow U=S\setminus T$
% slanted
\margin{\deleted{Autant dire}}

(\unsure{iv|iiv}) La topologie de $I$ est quotient final pour les applications $I'\mathbin{\text{---}}I$, $J\mathbin{\text{---}}I$ ($I'$ munie \deleted{celle de $I$, \unsure{conjointement}} \inserted{de $\overline{\tau'}$, $J$ \deleted{\unsure{de}} muni de $\overline{\tau}\mid J$}) (Ce qui exclut notamment la solution triviale pour (i) (ii), obtenue en prenant la topologie somme sur \deleted{$J'$} $I'=J'\underset{\text{disj.}}{\cup}\overset{\deleted{\sim}}{U'}$, sauf si $I$ lui$=$ est une somme top. $J\sqcup\underset{\substack{\Vert\\I\setminus J}}{U}$).
% slanted
\margin{C'est cette condition que j'avais d'abord oubliée, à la page 31 !}
% slanted
\margin{\deleted{Autant} dire que $\tau'$ sur $S$ est caractérisée par les propriétés analogues (i) (ii) (iii) (iv), la dernière \unsure{signifiant|signifie} que la top $\text{\textcircled{$\tau$}}$ sur $S$ est finale pour $\unsure{\tau''|\tau'}$ sur $S$ et $\tau$ sur $T$ (id $S\unsure{\Rightarrow|\rightrightarrows|\to}S$ $T\unsure{\hookrightarrow|\to}S$)}
% slanted
\margin{en fait, c'est indépendant aussi de l'hyp. \unsure{ultra-topologique}}

Démonstration

\deleted{a) et b).} \inserted{(a b c)} Il est clair que les ens. épinglés dans c) sont ceux qui satisfont $1^\circ)$, $2^\circ)$ $3^\circ)$ dans b), et que la représentation dans c) est unique. Il reste à voir, pour prouver (a,b,c), que les ens. forment une topologie, \deleted{et} qu'elle \deleted{\unsure{a les propriétés|a les propri.}} \inserted{est} dans $\underline{\mathscr{T}}$, et est moins fine des top $\in\underline{\mathscr{T}}$.

1) Ce sont les fermés d'une topologie.

Stables par intersections quelconques : clair

Stables par réunions finies : idem.

C'est $=$ stable par réunions quelconques s'il en est ainsi pour $(T,\tau'_0)$, $(\deleted{\unsure{I|D}}\unsure{S|I},\tau)$, d'où la première assertion de \unsure{d)|a)}.

\scanpage{44}\mspage{43}
Cette topologie est plus fine que $\tau\deleted{'}$, i.e. si $X$ est fermé dans $\tau$, il satisfait ces conditions : clair. Il induit $\tau'_0$, i.e. si $X\subset T$, il satisfait ces cond. sss il est fermé pour $\tau'_0$ : clair. Donc $\tau'\in\underline{\mathscr{T}}$. Prouvons que c'est un plus grand élément (top la moins fine). i.e que pour \deleted{tout} $\tau''\in\underline{\mathscr{T}}$, les fermés pour $\tau'$ sont des fermés pour $\tau''$. En effet, on a
\[
X=B\deleted{\cap}\cup\overline{X\capU}^{(\tau)}
\]
$B$ étant fermé dans $T$. (pour $\tau'_0$\inserted{\deleted{\unsure{donc}} pour $\tau^{\prime\prime\unsure{,}}\mid\overset{\sim}{T}$}), et $T$ étant \uline{fermé} dans $\deleted{\unsure{T|\tau}}S$ pour $\tau$ donc \inserted{a fortiori} pour $\tau''$, $B$ est fermé dans $S$ pour $\tau''$.
% slanted
\margin{(C'est la première fois qu'on utilise $T$ fermé dans $S$}
De $=$ $\overline{X\cap U}^{(\tau)}$ étant fermé pour $\tau$, l'est pour $\unsure{:|;}\tau''$, donc $X=B\cup\overline{X\cap U}^{(\tau)}$ l'est aussi. Cela prouve a) b) c).

Prouvons d)\inserted{(le début déjà prouvé)}, en nous rappelant que pour $x,y\in T$,
\[
\begin{aligned}
x\underset{\tau'}{\leq}y&\Longleftrightarrow\overline{x}^{(\tau')}\subset\overline{y}^{(\tau')}\\
&\Longleftrightarrow x\in\overline{y}^{(\tau')}.
\end{aligned}
\]
Comme $\tau'\mid U\ne\tau\mid U$, et $\tau'\mid T=\tau'_0$, les cas $x,y\in T$ et $x,y\in U$ sont déjà clairs. Reste le cas où soit $x\in U,y\in T$, soit l'inverse. Le premier cas ne peut se produire, car si $y\in T$, $\overline{y}^{(\tau')}\inT$ ($T$ \uline{fermé} dans $S$ pour $\tau'$) donc $x\notin\overline{y}^{(\tau')}$ \deleted{puisque} si $x\in U$. Donc reste le cas $[x\in T,]$ $y\in U$, prouver $x\underset{\tau'}{\leq}y\Longleftrightarrow x\underset{\tau'}{\leq}y$
\margin{NB c'est vrai sans supposer $x\inT$}

